\begin{abstract}
The differential emission measure (DEM) records how much plasma lies at each temperature along a line of sight through the solar corona. Classical solvers recover it from extreme-ultraviolet images one pixel at a time, at a cost of about four million constrained optimizations per image, and the supervised networks that make this fast must still be trained on solver output. We instead train a per-pixel network on the solver's own objective, so no solver DEM is ever a training target. On 153 held-out images, scored on the same pixels as a supervised emulator, our Elastic Net model leads on five of six headline measures and ties the sixth. Our Basis Pursuit model leads on three of six in all three training seeds, with a median pixel error 3.6 times lower, while the supervised model keeps lower DEM MSE (0.91 against 1.08 to 1.35) and AIA MAE (2.95 against 3.54). Basis Pursuit has a hard noise band, which we replace with a squared penalty, and the penalty weight then decides what the network can learn. Above a threshold computable from the linear program, each channel settles outside its band by an amount set by its dual price and noise, falling as $1/\mu$, and at 10,899 held-out pixels the numerical minimizer of the loss matches this prediction (median ratio 1.00). Below the threshold the solution changes support, and at lower weights hot flare plasma leaves it altogether. This explains a failure at flare cores that more training could not fix, and it lets us choose the weight from the instrument response and the noise alone.
\end{abstract}
